\section{Methodology}
\label{sec:methodology}

This section describes the research design, data sources, mathematical formulation, and implementation details of the proposed carbon-intelligent workload scheduling model.

\subsection{Research Design}

This thesis applies a \emph{deterministic optimization} approach to carbon-aware scheduling.
The core method is a linear program (LP) that, given historical hourly carbon intensity (CI) and renewable fraction (RF) data, computes the globally optimal assignment of compute load across five regions and 17,544 hours.
Optimality guarantees are evaluated through rolling-window backtesting against three baselines: carbon-agnostic scheduling, greedy temporal shifting, and an oracle with perfect foresight.
This design falls within the \emph{validation of methodology} category: it applies and extends established LP-based scheduling frameworks~\cite{liu2012, riepin2024spatiotemporal} to a new dataset and region portfolio, with a novel RF weighting term as the primary methodological contribution.

\subsection{System Model}

We consider a cloud service operator managing \emph{delay-tolerant} AI workloads across $|\mathcal{R}| = 5$ geographically distributed datacenter regions.
Each region $r \in \mathcal{R}$ is connected to a distinct electricity grid.
The planning horizon spans $|\mathcal{T}| = 17{,}544$ hourly time steps (2024-01-01 to 2025-12-31 UTC).

\textbf{Workload model.}
Total compute demand $D$ (kWh) must be satisfied over the horizon.
At each hour $t$, region $r$ can accept up to $C_{\max,r}$ kWh.
A fraction of load is \emph{flexible}: it can be shifted in time or routed to another region.
Latency-sensitive (inflexible) load is excluded.

\textbf{Grid regions.}
Regions are selected to represent structurally distinct grid types:

\begin{table}[h]
\centering
\begin{tabularx}{\linewidth}{p{2.2cm}p{2.2cm}p{3.8cm}X}
\toprule
\textbf{Region} & \textbf{Zone ID} & \textbf{Grid character} & \textbf{CI profile} \\
\midrule
PJM       & US-MIDA-PJM & Coal/gas/nuclear mix; high variability & High, volatile \\
NYISO     & US-NY-NYIS  & Hydro + nuclear; some gas peaking      & Medium, diurnal \\
Finland   & FI          & Hydro + nuclear + wind; Nordic grid    & Low, seasonal \\
Belgium   & BE          & High nuclear; European interconnect    & Medium-low \\
Singapore & SG          & Near-total natural gas; isolated grid  & High, stable \\
\bottomrule
\end{tabularx}
\caption{Five target regions and their grid characteristics. Singapore serves as a control case: low CI variability implies limited temporal shifting benefit.}
\label{tab:regions}
\end{table}

\subsection{Data Sources}

\textbf{Carbon intensity (CI).}
Hourly CI values (gCO\textsubscript{2}eq/kWh) are retrieved from the ElectricityMaps API via the \texttt{/v3/carbon-intensity/past} endpoint.
ElectricityMaps aggregates raw TSO measurements and applies ML models to produce a continuous, pre-validated series~\cite{radovanovic2023}.
Coverage: 2024-01-01 00:00 to 2025-12-31 23:00 UTC; 17,544 observations per region; zero missing values.

\textbf{Renewable fraction (RF).}
Computed from the \texttt{/v3/power-breakdown/past} endpoint:
\[
  \mathrm{RF}_r(t) = \frac{\sum_{s \in \mathcal{S}_{\mathrm{clean}}} P_{r,s}(t)}{\sum_{s} P_{r,s}(t)},
  \quad \mathcal{S}_{\mathrm{clean}} = \{\text{wind, solar, hydro, geothermal, biomass}\}
\]
Nuclear is excluded: it is low-carbon but dispatchable and non-intermittent, making it a distinct scheduling consideration.

\textbf{Electricity price (budget parameter).}
Per-region price series parameterize Constraint C3.
Sources: PJM Locational Marginal Prices, NYISO day-ahead prices, ENTSO-E (Finland, Belgium), EMC spot prices (Singapore).
\todo{Confirm data availability and add source citations.}

\subsection{Key Indicators}

\begin{itemize}[leftmargin=1.5em]
  \item \textbf{CI$_r(t)$}: carbon cost of one kWh of compute in region $r$ at hour $t$. Primary optimization signal.
  \item \textbf{RF$_r(t)$}: renewable fraction at time $t$. Enters the objective to incentivize variable renewable absorption.
  \item \textbf{CV$_r$}: coefficient of variation of CI, $\sigma(\mathrm{CI}_r)/\mu(\mathrm{CI}_r)$. Computed post-hoc from the full 2024--2025 historical series; not an LP input. Used as a diagnostic to explain cross-regional differences in achievable savings after backtesting.
\end{itemize}

\subsection{Mathematical Formulation}

\subsubsection*{Decision Variables}

Let $x_{r,t} \geq 0$ denote the compute load (kWh) assigned to region $r$ at hour $t$.
This single variable encodes both temporal scheduling and spatial routing simultaneously, enabling globally optimal joint allocation without a sequential two-phase decomposition.

\subsubsection*{Objective Function}

\begin{equation}
  \min_{x} \quad \sum_{r \in \mathcal{R}} \sum_{t \in \mathcal{T}}
    x_{r,t} \cdot \mathrm{CI}_r(t) \cdot \bigl(1 - \alpha \cdot \mathrm{RF}_r(t)\bigr)
  \label{eq:objective}
\end{equation}

CI and RF are both deterministic input parameters, not decision variables. $\mathrm{CI}_r(t)$ is the direct carbon outcome: actual emissions per kWh consumed, and the primary optimization signal~\cite{radovanovic2023}. $\mathrm{RF}_r(t)$ is a supply-structure signal: the share of variable renewable generation. Two hours with the same CI can differ in RF — one driven by wind surplus, one by nuclear — and RF discriminates between them~\cite{radovanovic2023, liu2012}.

The multiplicative discount form $(1 - \alpha \cdot \mathrm{RF}_r(t))$ is derived from OLS regression of CI on RF across all five regions (EDA, Section~10). The normalized slope $\beta_1/\overline{\mathrm{CI}}$ is more consistent across regions than the raw slope $\beta_1$, indicating that the RF effect scales proportionally with the ambient CI level rather than acting as a fixed absolute reduction. This supports the multiplicative form over an additive alternative. Parameter $\alpha \in [0, 1]$ is estimated as $\approx -\beta_1/\beta_0$ from the regression; a single shared value is used if the per-region estimates cluster closely, otherwise region-specific values apply. For any region where $|r(\mathrm{CI}, \mathrm{RF})| < 0.2$, $\alpha$ is set to zero.

\subsubsection*{Constraints}

\begin{align}
  &\sum_{r \in \mathcal{R}} \sum_{t \in \mathcal{T}} x_{r,t} = D
  \tag{C1} \label{eq:c1} \\[4pt]
  &x_{r,t} \leq C_{\max,r} \quad \forall\, r,\; t
  \tag{C2} \label{eq:c2} \\[4pt]
  &\sum_{t \in \mathcal{T}} x_{r,t} \cdot P_r(t) \leq B_r \quad \forall\, r
  \tag{C3} \label{eq:c3} \\[4pt]
  &x_{r,t} \geq 0 \quad \forall\, r,\; t
  \tag{C4} \label{eq:c4}
\end{align}

C1 enforces exact demand satisfaction~\cite{liu2012}.
C2 bounds regional server capacity~\cite{liu2012}.
C3 caps per-region electricity expenditure at budget $B_r$, with prices in local currency to avoid cross-region normalization~\cite{hao2024joint}. Electricity price data collection is in progress; C3 will be incorporated once sources for all five regions are confirmed.
C4 is non-negativity.
The current implementation uses C1, C2, and C4 as the operative constraints. C3 will be added once electricity price data is confirmed.

The supervisor's VCC mechanism~\cite{colangelo2025} — which bounds the flexible component of load against a CI-indexed capacity ceiling — is a planned extension. It will be incorporated as an additional constraint once Google Cluster Trace data is available to generate the VCC series.

\begin{table}[h]
\centering
\begin{tabularx}{\linewidth}{p{3.6cm}Xp{4cm}}
\toprule
\textbf{Component} & \textbf{Description} & \textbf{Reference} \\
\midrule
Objective \eqref{eq:objective}     & CI-weighted carbon cost signal          & \citet{radovanovic2023} \\
RF weighting $(1-\alpha\cdot\mathrm{RF})$ & Multiplicative discount derived from EDA regression; motivation from \citet{radovanovic2023}, \citet{liu2012}; form original to this thesis & This thesis \\
C1--C2                             & Demand and capacity bounds              & \citet{liu2012} \\
C3 (budget)                        & Per-region electricity cost cap (pending) & \citet{hao2024joint} \\
VCC (planned)                      & Flexible load ceiling, pending data     & \citet{colangelo2025} \\
Joint LP structure                 & Spatio-temporal allocation              & \citet{riepin2024spatiotemporal}; \citet{attenni2024shifting} \\
\bottomrule
\end{tabularx}
\caption{Component-level references for the mathematical formulation.}
\label{tab:modelrefs}
\end{table}

\subsection{Implementation}

\textbf{Solver.}
The model is implemented in Python using the \texttt{gurobipy} interface to Gurobi 10.
With $|\mathcal{R}|=5$ and $|\mathcal{T}|=17{,}544$, the LP has 87,720 continuous variables and approximately 105,264 constraints.
\todo{Report actual solve time after implementation.}

\textbf{Rolling-window backtesting.}
The LP is solved in a rolling-window fashion with decision horizon $W$ and look-ahead $L$ to simulate real-world scheduling.
\todo{Specify W and L; report sensitivity to window size.}

\textbf{Data pipeline.}
CI and RF inputs are loaded from pre-computed Parquet files.
\todo{Electricity price source: specify and cite.}
All inputs are aligned to hourly UTC timestamps before passing to the solver.
Code and data are available at \url{https://github.com/NOSIEMPRE/AI-Carbon-Optimization}.

\subsection{Heuristic Fallback}

If rolling-window LP solve times are prohibitive, a greedy carbon-aware heuristic serves as fallback and as an additional baseline:

\begin{algorithm}
\caption{Greedy Carbon-Aware Shifting}
\begin{algorithmic}[1]
\REQUIRE Remaining demand $d$, window $[t_1, t_2]$, regions $\mathcal{R}$, CI, RF, capacities
\STATE Compute $\hat{c}_{r,t} = \mathrm{CI}_r(t)(1 - \alpha \cdot \mathrm{RF}_r(t))$ for all $(r,t)$
\STATE Sort $(r,t)$ pairs by $\hat{c}_{r,t}$ ascending
\FOR{each $(r,t)$ in sorted order}
  \STATE $x_{r,t} \leftarrow \min(d,\; C_{\max,r})$;\quad $d \leftarrow d - x_{r,t}$
  \IF{$d = 0$} \STATE \textbf{break} \ENDIF
\ENDFOR
\end{algorithmic}
\end{algorithm}

\subsection{Assumptions and Scope}

\begin{enumerate}[leftmargin=1.5em]
  \item \textbf{Perfect historical information.} The LP uses the full 2024--2025 CI/RF series; results represent an upper bound on real-time savings.
  \item \textbf{Homogeneous compute-to-power conversion.} PUE differences across regions are abstracted away.
  \item \textbf{No cross-region data transfer costs.} Network latency is not modeled; the setup is appropriate for batch workloads.
  \item \textbf{Stationary demand.} Total demand $D$ is fixed; load shedding is not considered.
\end{enumerate}
